% !PRINTED none
% Abschnitt-opening page: no running head / page number visible in the
% scanned bands. Upper part of the page is blank, followed by the
% centred title block, then §187 begins with running text; a folio may
% exist elsewhere on the physical page but is not visible here.

\begin{center}
\gesperrt{Sechsundzwanzigster Abschnitt.}

\textbf{Zahlenwerte der algebraischen Funktionen.}

\rule{4cm}{0.4pt}

\textbf{§ 187. \gesperrt{Der Punkt.}}
\end{center}

Es entsteht nun die Frage, wie man den algebraischen Funktionen
 eine Bedeutung im Gebiete der \gesperrt{Zahlen} beilegen kann. Wir
setzen dabei das Gebiet der reellen und imaginären rationalen
und irrationalen Zahlen als gegeben voraus, und in diesem
Gebiete das Rechnen mit den vier Spezies. Wir erweitern dieses
Gebiet durch Hinzufügung eines Zahlzeichens \glqq Unendlich\grqq{}: $\infty$,
und rechnen auch mit diesem Zeichen in bekannter Weise, so daß

\weq{1}{\frac{1}{\infty} = 0, \quad \frac{1}{0} = \infty, \quad \infty \cdot \infty = \infty}

ist, während den Symbolen
\[
\infty \pm \infty, \quad 0 \cdot \infty, \quad \frac{0}{0}, \quad \frac{\infty}{\infty}
\]
\gesperrt{keine Bedeutung zukommt} (oder auch, in jedem einzelnen
Falle nach Bedarf, ein beliebiger Zahlwert beigelegt wird).

Der Körper $\Omega$ besteht jetzt aus allen möglichen algebraischen
Funktionen $\alpha, \beta, \gamma, \dots$, die nach § 186 als rationale Funktionen
von zweien unter ihnen mit Zahlenkoeffizienten dargestellt werden
können. Auf die Art dieser Darstellung kommt es zunächst
nicht an.

Wir stellen nun folgende Definition auf:

1. \gesperrt{Wenn alle Individuen $\alpha, \beta, \gamma, \dots$ des Körpers $\Omega$
durch bestimmte Zahlwerte $\alpha_0, \beta_0, \gamma_0, \dots$ in der Weise
ersetzt werden, daß}

\begin{center}
\begin{minipage}{0.7\linewidth}
1. \quad $\alpha = \alpha_0$, wenn $\alpha$ konstant,

2. \quad $(\alpha + \beta)_0 = \alpha_0 + \beta_0$,

3. \quad $(\alpha - \beta)_0 = \alpha_0 - \beta_0$,

4. \quad $(\alpha\beta)_0 = \alpha_0\beta_0$,

5. \quad $\left(\dfrac{\alpha}{\beta}\right)_0 = \dfrac{\alpha_0}{\beta_0}$,
\end{minipage}
\end{center}
%CONT>>
